\begin{textsample}{POS Dim 1 – se – Score 29.00 – se/se\_em\_7.txt}  \label{ex:f1_pos_119}
3.1.4 Língua Inglesa Nas palavras de Paulo Freire, “ saber \textbf{que} ensinar não é transferir conhecimento, mas criar as possibilidades para a sua própria produção ou a sua construção ” ( \textbf{FREIRE}, 2018 ) é \textbf{basilar} para \textbf{um} currículo \textbf{que} transmite visões sociais e \textbf{que} produz identidades individuais numa sociedade ampla e tão complexa.
Assim, o referencial curricular sergipano do Ensino Médio, no \textbf{que} \textbf{diz} respeito à língua inglesa, foi construído visando contribuir na formação dos jovens, proporcionando a sua inserção no contexto socioeducacional e ampliando a visão de mundo deles para os aspectos interculturais da língua inglesa.
Nesse sentido, a instituição educacional, \textbf{enquanto} espaço de encontro sociocultural, projeta a vida social marcada pela pluralidade de \textbf{sujeitos}, pensamentos, valores, atitudes e crenças, o \textbf{que} \textbf{remete} a \textbf{uma} confluência entre língua e cultura \textbf{que}, segundo Silva ( 2019 ), quando direcionada por \textbf{uma} \textbf{ótica} crítica, viabiliza o fomento de \textbf{uma} educação intercultural \textbf{que} transcende a \textbf{mera} descoberta do outro ou o (re)conhecimento daquilo \textbf{que} não se assemelha aos seus hábitos e suas vivências.
Ainda de acordo com Silva apud Siqueira \& Sousa ( 2019 ), para dar conta da realidade cada vez mais plural da língua inglesa, o professor precisa orientar -se por “ \textbf{uma} \textbf{abordagem} de base intercultural, fundada em atitudes democráticas e de acolhimento às diferenças ”, o \textbf{que} leva a diferentes \textbf{sujeitos} interagirem com base em práticas \textbf{dialógicas} \textbf{que} explorem e valorizem a diversidade relacionada ao ambiente da instituição educacional.
Dessa forma, a Base Nacional Comum Curricular ( BNCC ) propõe \textbf{que} os estudantes tenham a possibilidade de explorar a presença da multiplicidade de usos da língua inglesa na cultura digital, bem como nas culturas juvenis, além da possibilidade de apropriar -se de estudos e pesquisas acadêmicas.
Nessa perspectiva, entende -se \textbf{que} o caráter intercultural da língua é essencial para a formação integral do estudante, \textbf{uma} vez \textbf{que} o indivíduo demonstra a sua cultura pelo uso \textbf{que} se faz da língua.
Segundo Melo ( 2016 ) apud Hall ( 2012 ), é \textbf{preciso} \textbf{ver} a variedade cultural dos estudantes não como \textbf{uma} versão deficiente de \textbf{uma} noção idealizada de língua e cultura, mas como ferramenta \textbf{que} permite sua participação em diversos grupos e comunidades em seu mundo sociocultural.
O ensino da língua inglesa não se \textbf{restringe} apenas aos aspectos estruturais da língua, engloba todo o universo cultural no qual está inserido.
Isso \textbf{implica} no desenvolvimento de competências e habilidades específicas para apropriação do idioma ; assim, as aprendizagens essenciais para o desenvolvimento das habilidades e competências em inglês, como propõe a BNCC, culminarão com a amplitude da capacidade discursiva do \textbf{discente} e de vivências experimentais com diversas práticas de linguagem, de maneira \textbf{que} possibilite aos jovens desenvolver habilidades de comunicação e comportamento diferenciados.
Reforça -se \textbf{aqui} \textbf{que} o Currículo de Sergipe Etapa Ensino Médio está \textbf{embasado} sob o viés dos documentos normativos, tais como : Base Nacional Comum Curricular ( 2018 ), as Diretrizes Curriculares Nacionais para o Ensino Médio, Lei 13.415/17, Lei 9.394/96 ( Lei das Diretrizes e Bases da Educação Nacional ), Portaria nº 1.432/18, dentre outros ; \textbf{que} apresentam o ensino da língua inglesa através de \textbf{um} caráter intercultural para atender às necessidades dos jovens em seu processo de formação.
Com efeito, a Lei 13.415/17 traz, em seu art. 3º, § 4º, \textbf{que} os currículos do ensino médio incluirão, obrigatoriamente, o estudo da língua inglesa e poderão ofertar outras línguas estrangeiras, em caráter optativo, apontando, preferencialmente, o espanhol, desde de \textbf{que} haja disponibilidade de oferta de vagas, com locais e horários definidos pelos sistemas de ensino.
Desse modo, a necessidade de mudança é \textbf{um} fator primordial, tanto nos currículos como nas formas de ensino e aprendizagem.
A BNCC nos traz competências específicas da língua \textbf{que} contribuem para a formação de estudantes críticos, responsáveis, solidários, \textbf{que} serão \textbf{aprendizes} por toda a vida.
Sendo assim, o Currículo de Sergipe Etapa Ensino Médio vem assegurar \textbf{que} os estudantes desenvolvam essas competências, permitindo -lhes \textbf{que} se comuniquem em língua inglesa, possibilitando o alcance de seu projeto de vida bem como a sua inclusão no mundo do trabalho.
Com base nesse entendimento, o ensino de língua inglesa, por meio de \textbf{um} letramento crítico, leva o indivíduo a ter acesso ao conhecimento, a compreender outras culturas e outras visões de mundo, além de desenvolver as potencialidades \textbf{que} os coloquem como \textbf{atores} de seu próprio desenvolvimento.
Assim, Piorino ( 2011 ) afirma : > É \textbf{urgente} \textbf{que} tenhamos : [ … ] \textbf{um} currículo escolar mais aberto, flexível e investigativo, construído a partir das necessidades daqueles \textbf{que} o vivenciam, mobilizando interesses de alunos e professores, de modo \textbf{que} estes se reconheçam e \textbf{vejam} sentido nas \textbf{tarefas} desenvolvidas com ou sem tecnologias.
De igual modo, no \textbf{que} concerne ao ensino do inglês para o alcance de \textbf{uma} educação integral e para atender à proposta curricular, a língua sempre esteve presente como \textbf{um} recurso importante para se ter acesso a bens culturais e científicos produzidos em outros contextos sociais bem como em outros espaços geográficos, e como \textbf{desdobramento}, o processo de internacionalização e o desenvolvimento das tecnologias digitais da informação e da comunicação têm transformado o papel da língua inglesa, contribuindo para o surgimento de novas formas de conhecer e de produzir conhecimento.
Esse processo de ensino e aprendizagem de língua inglesa, \textbf{que} perpassa pela internacionalização e pelo desenvolvimento das tecnologias digitais, inicia -se no Ensino Fundamental anos finais e já vem revestido da importância do contexto cultural, para \textbf{que} dessa forma \textbf{possamos} ter \textbf{uma} prática efetiva, pois \textbf{implica} em \textbf{um} novo papel \textbf{que} a língua vem a exercer, o de língua franca.
Nesse viés, o ensino de língua estrangeira \textbf{moderna} no Ensino Médio objetiva auxiliar os estudantes a compreender como as línguas contribuem para a sua constituição como \textbf{sujeitos}, através de valores construídos nas práticas sociais, ajudando -os em sua formação para a cidadania.
Afinal, os nossos estudantes são, segundo Jordão ( 2011 ), agentes, fazedores.
Eles buscam, selecionam, aprendem e interagem.
Nesse contexto, o inglês assume \textbf{um} papel fundamental na transformação social do indivíduo e, para tanto, o seu ensino deve ser direcionado por \textbf{uma} perspectiva crítica, contribuindo para participação ativa do \textbf{discente}.
Portanto, entende -se \textbf{que} para o alcance de \textbf{uma} educação com foco na formação integral do indivíduo, os estudantes devem estar no centro do processo de ensino e aprendizagem, tendo a oportunidade de trabalhar o autoconhecimento, o pensamento crítico, a criatividade, a sociabilidade, através de \textbf{um} ensino personalizado, \textbf{que} contribui na construção \textbf{identitária} desses jovens, no agir e posicionar -se criticamente na sociedade, na ampliação de repertório linguístico, priorizando os aspectos de inteligibilidade e compartilhamento de informações.

% matched lemmas: abordagem, aprendiz, aqui, ator, basilar, desdobramento, dialógico, discente, dizer, embasar, enquanto, freire, identitária, implicar, mero, moderno, possar, preciso, que, remeter, restringir, sujeito, tarefa, um, urgente, ver, ótica
\end{textsample}
